\scanpage{30}
\begingroup\csname tagsleft@true\endcsname

% top, first member
or le premier membre (grâce à (2.3.2)) est $\phi^{n'}(x,my)=m\,\phi^{n'}(x,y)$. La compatibilité \deleted{(2.3.1)} (2.\deleted{2}3.2) \deleted{est}\inserted{implique} à son tour \deleted{conséquence de} la compatibilité symétrique\inserted{ de (2.3.1)}, exprimée par la commutativité de

% upper square, label at left
\diagram{\begin{tikzcd}[ampersand replacement=\&,column sep=small]
\& \deleted{\unsure{n'|m|\mathrm{id}}}\;{}_nP\arrow[d,"\handwritten{\mathrm{id}}"] \& \mathsf{x} \& {}_{n'}Q\arrow[d,"\handwritten{m}"]\arrow[rr,"{\phi_{\handwritten{\xi}}^{n'}\handwritten{|}\,{}_nP\;\mathsf{x}\;{}_{\handwritten{n'}}Q}"] \& \& {}_nG\arrow[d,"\handwritten{\mathrm{id}}"]\\
\text{(2.3.\deleted{\unsure{3}}4)} \& {}_nP \& \mathsf{x} \& {}_n\handwritten{Q}\arrow[rr,"{\phi_{\handwritten{\xi}}^n}"] \& \& {}_nG\,, 
\end{tikzcd}}

% below square
qui s'écrit en formules

% upper formula
\[\phi^n(x,my)=\phi^{n'}(x,y)\qquad (x\in{}_nP(S),\ y\in{}_nQ(S)).\tag*{(2.3.\deleted{\unsure{4}}5)}\]

% middle, crossed opening and reduction
\deleted{Montrons comment (2.3.\unsure{2|4}) implique (2.3.\unsure{4|1|3}) : avec les notations de cette dernière formule,} Comme $\,{}_{mn'}P\handwritten{\xrightarrow{\mathrm m}}{}_nP$ est un épimorphisme, on est ramené à prouver (2.3.2) dans le cas où $x=mx'$, où $x'\in{}_{n'}P(S)$, \deleted{donc}\inserted{or} on a $\deleted{\phi^n(mx,y)=\phi^n(m^{\inserted{\handwritten{2}}}x',y)=\phi^n}\quad\phi_{\handwritten{\xi}}^n(mx',my)\overset{\inserted{\text{(2.3.3)}}}{=}m\,\phi_{\handwritten{\xi}}^{n\handwritten{'}}(x',y)=\phi_{\handwritten{\xi}}^{n'}(mx',y)$ (par bilinéarité de $\phi^{n'}$) $=\phi_{\handwritten{\xi}}^{n'}(x,y)$ puisque $x=mx'$ \deleted{\unsure{q|g}}dans $\,{}_{n'}P(S)$, ce qui établit (2.3.5).

% lower middle, functoriality
Prouvons donc la compatibilité (2.3.1). La fonctorialité\inserted{ en $Q$} de l'homomorphisme (2.1.2) nous montre qu'on peut supposer $nQ=0$. Alors (2.1.2) est le composé

% lower middle, composite
\[\mathrm{Biext}^1(P,Q;G)\handwritten{\xrightarrow{\sim}}\mathrm{Hom}(\underline{\mathrm{Tor}}_1(P,Q),G)\handwritten{\xleftarrow{\sim}}\mathrm{Hom}({}_nP\otimes Q,G)\,,\]

% lower middle, cancelled line
\deleted{en $Q$ nous montre que (2.1.2) est le composé}

% lower middle, cancelled display
\[\deleted{\mathrm{Biext}^1(P,Q;G)\qquad\mathrm{Hom}(\mathrm{Tor}_1(P,Q),G)\qquad\mathrm{Hom}(\mathrm{Tor}_1(P,{}_nQ),G)}\]

% lower, two arrows
où la première flèche est celle de (1.1.3) et ne dépend pas du choix de $n$, et où la deuxième flèche est transposée de l'isomorphisme $\underline{\mathrm{Tor}}_1(P,Q)\handwritten{\xrightarrow{\overset{\unsure{\partial|d}}{\sim}}}{}_nP\otimes Q$ défini dans (1.3.4). Donc la compatibilité voulue équivaut essentiellement à \deleted{celle} la commutativité du diagramme

% foot, triangle
\diagram{\begin{tikzcd}[ampersand replacement=\&]
\& {}_{n'}P\otimes Q\arrow[dd,"{\handwritten{m\otimes \mathrm{id}_Q}}"]\\
\underline{\mathrm{Tor}}_1(P,Q)\arrow[ur]\arrow[dr] \&\\
\& {}_nP\otimes Q\quad\deleted{\handwritten{\unsure{m|+}}}\quad\deleted{\handwritten{\unsure{q|4}}}\,,
\end{tikzcd}}

% foot, last lines
où les flèches obliques sont les deux isomorphismes du type (1.3.4), relatifs à $n$ et $n'$ respectivement. Or cette compatibilité résulte de
\endgroup
% archive sheet number: 29
